\chapter{Research Methodology}
\label{ch:methodology}

\section{Research Methodology}
The study is a retrospective, observational modelling study on de-identified
critical-care data, conducted as quantitative empirical research. The research
question (Section~\ref{sec:rq}) is operationalised as a prediction task with
pre-registered baselines (SIRS, qSOFA, SOFA) and an explicit null hypothesis for
the proportional-hazards assumption. A deliberate design choice constrains all
modelling to classical statistics; this preserves interpretability, supports
formal assumption testing, and establishes a transparent, transportable baseline
against which machine-learning approaches can later be compared on identical
data splits.

The methodology follows the standard prognostic-model development cycle: cohort
definition, outcome labelling, feature construction, model fitting, internal
validation, and external validation. Two principles govern every stage. The
first is \emph{leakage safety}: at any prediction time only information that
would genuinely be available to a clinician at that hour is used, implemented
through forward-filled values and time-since-last-measurement counters. The
second is \emph{honest evaluation}: discrimination is measured out-of-sample
using patient-grouped cross-validation, and the frozen model is transported to a
second, independent database.

\section{Software Development Methodology}
The pipeline was developed incrementally as a sequence of numbered, single-purpose
stages (Chapter~\ref{ch:design}), each writing an inspectable intermediate
artefact that the next stage consumes. This staged design supports
reproducibility and debuggability: any stage can be re-run in isolation, and a
single configuration module is the sole source of truth for file paths, clinical
item identifiers, and thresholds. Data extraction and panel construction are
implemented in Python (pandas); statistical modelling and reporting are
implemented in R, with each analysis tier emitting a self-contained report. All
stochastic steps use a fixed random seed, and the multi-gigabyte source tables
are processed by streaming in bounded-memory chunks rather than being loaded
whole.

\section{Data}
The development data is MIMIC-IV v3.1 \cite{johnson2023mimiciv}, a large
single-centre critical-care database. External validation uses the eICU
Collaborative Research Database v2.0 \cite{pollard2018eicu}, a multi-centre
collection spanning more than 200 United States ICUs. Both are de-identified and
were accessed under the required credentialing.

\section{Cohort and Outcome Labelling}
An adult ICU cohort was constructed (age $\ge 18$ years, first ICU stay only,
length of stay $\ge 4$\,h), yielding 65{,}241 stays. Suspected infection follows
Seymour \emph{et al.} \cite{seymour2016criteria}: a qualifying antibiotic order
paired with a body-fluid culture within the Sepsis-3 time windows, with
surveillance swabs (for example, MRSA screens) excluded to avoid inflating the
suspicion rate. Sepsis-3 is defined as suspected infection together with an
acute rise in the Sequential Organ Failure Assessment (SOFA) score of at least
two points; the cohort prevalence is 37.1\%, with septic shock at 9.2\%, and
in-hospital mortality of 18.7\% in the septic group versus 6.1\% in the
non-septic group, consistent with a valid label.

\section{Hourly Panel and Onset Definition}
Every stay was resampled to a regular hourly grid (approximately 3.1 million
stay-hours). For each variable, three quantities are retained per hour: the raw
value, a forward-filled value (what the clinician currently knows), and a
time-since-last-measured counter that encodes irregular sampling as an explicit
signal. Hourly SOFA is scored from the forward-filled values, the exact onset
hour is located as the first hour meeting the Sepsis-3 rise inside the suspicion
window, and a per-hour ``onset within 12\,h'' label is derived for early-warning
evaluation. An hour-six landmark yields 48{,}829 at-risk stays (those that
reached hour six without already being septic), which is the target population
for early prediction.

\section{Models}
The baselines are the rule-based scores SIRS, qSOFA, and NEWS computed directly
from vitals. The statistical models comprise: (i) a static Cox model with
Schoenfeld testing and a time-varying reformulation on failure; (ii) an
elastic-net penalised logistic \emph{nomogram} at the hour-six landmark using an
informative-missingness design (each laboratory value is median-imputed and
accompanied by a binary ``was measured'' indicator); (iii) a \emph{landmark
supermodel} stacked across landmarks at 6, 12, 18, 24, 36, and 48\,h so that one
model emits hourly-updating risk; (iv) generalised estimating equations (GEE)
with cluster-robust standard errors for within-patient correlation; (v)
group-based trajectory modelling of first-24\,h SOFA trajectories linked to
onset and mortality; and (vi) confounder-adjusted causal estimates (propensity-score
matching, inverse-probability weighting, augmented IPW, and doubly robust
estimation) for early antibiotics.

\section{Evaluation and External Validation}
Discrimination uses the area under the receiver operating characteristic curve
(AUROC) with DeLong confidence intervals; calibration uses the Brier score,
calibration-in-the-large, and the calibration slope; clinical utility uses the
PhysioNet/CinC 2019 utility score \cite{reyna2020physionet} and an alarm-burden
curve (false alarms per patient-day at fixed sensitivity). For external
validation, the hour-six landmark table is rebuilt on eICU using the identical
cohort, panel, SOFA, onset, and feature definitions; the frozen MIMIC
coefficients are then applied without refitting, and the frozen-transport
performance is compared against an eICU-native refit that establishes the best
discrimination the same features can achieve in that population.
